% Generated draft: requires mathematical and source review.
\documentclass[11pt]{article}
\usepackage[margin=25mm]{geometry}
\usepackage{fontspec}
\setmainfont{Latin Modern Roman}
\usepackage{amsmath,amssymb,mathtools,unicode-math}
\setmathfont{Latin Modern Math}
\usepackage{enumitem,needspace,xurl,hyperref}
\setlength{\parindent}{0pt}
\setlength{\parskip}{5pt}
\setlength{\emergencystretch}{3em}
\newcommand{\N}{\mathbb{N}}
\newcommand{\Z}{\mathbb{Z}}
\newcommand{\Q}{\mathbb{Q}}
\newcommand{\R}{\mathbb{R}}
\newcommand{\C}{\mathbb{C}}
\newcommand{\F}{\mathbb{F}}
\newcommand{\A}{\mathbb{A}}
\newcommand{\PP}{\mathbb{P}}
\newcommand{\E}{\mathbb{E}}
\newcommand{\eps}{\varepsilon}
\newcommand{\dd}{\,\mathrm{d}}
\begin{document}
\section*{UnsolvedMath ID: 30004585}
\textbf{Source problem code:} catalogue entry
\subsection{Definitions and mathematical statement}
To explain the stated equivalence, use the usual convention of one-sorted finitary universal algebra with nonempty underlying sets. This convention is an interpretation: the supplied record does not specify the permitted signatures or algebraic setting.

Fix an algebraic signature \(L\). An \(L\)-identity is a universally quantified equation \(s\approx t\) between \(L\)-terms. A variety is a class \(\mathcal V=\operatorname{Mod}_L(T)\) of all \(L\)-algebras satisfying a set \(T\) of identities. Write \(\mathcal T_L\) for the trivial variety consisting of all one-element \(L\)-algebras. A variety \(\mathcal V\) is \emph{minimal} if
\[
\mathcal V\ne\mathcal T_L
\quad\text{and}\quad
\forall\mathcal W\,
\bigl(\mathcal W\text{ is an }L\text{-subvariety of }\mathcal V
\Longrightarrow
\mathcal W=\mathcal T_L\text{ or }\mathcal W=\mathcal V\bigr).
\]
An equational theory \(T\) is a set of identities closed under equational consequence. It is \emph{equationally complete} in the corresponding nontrivial sense if it has a model with at least two elements and, for every \(L\)-identity \(s\approx t\),
\[
T\not\models s\approx t
\quad\Longrightarrow\quad
T\cup\{s\approx t\}\models x\approx y,
\]
where \(x,y\) are distinct variables and all variables in identities are universally quantified. The identity \(x\approx y\) forces every nonempty model to have one element. Under these conventions, an equationally complete theory is exactly the equational theory of a minimal variety.

\textbf{Source question.} Can equationally complete theories, equivalently minimal varieties, be classified?

The requested task is a general classification. The record does not specify whether the signature is fixed or variable, whether restrictions such as a finite signature are intended, the equivalence relation for classification, or the required form of a classification. These remain unresolved qualifications; the definitions above explain a conventional interpretation without supplying those missing specifications.
\subsection{Short English statement}
The question asks for a classification of nontrivial varieties whose only proper subvariety is the trivial variety. Equivalently, it asks for a classification of equational theories for which adding any identity not already implied forces all models to have one element. The permitted signatures and the meaning of classification are unspecified. The supplied record labels the problem open; this draft does not independently verify that status.
\subsection{Sources}
\begin{itemize}
\item Local UnsolvedMath registry, supplied complete record: permanent ID \(30004585\), title \emph{Classification of Equationally Complete Theories}, exact source problem code \texttt{catalogue entry}. The original, clean, and displayed statements coincide. The record marks the statement as exact and publication-ready; those are supplied editorial assessments. No individual proposer is supplied.
\item Supplied source citation: \emph{Homogeneous Structures: Model Theory meets Universal Algebra} (2022), Oberwolfach Reports, DOI 10.4171/owr/2021/1, pp.\ 5--32. Source URL and supplied literature-source link, labeled ``Oberwolfach report'': \url{https://doi.org/10.4171/owr/2021/1}. The record gives proposed year 2021; the relationship between that year and the citation's 2022 date is not explained.
\item Supplied literature assessment, checked \texttt{2026-08-22}: the record labels the problem open and states that the report documents substantial partial classifications of minimal varieties, including important abelian cases, while the requested general classification remains unresolved. This is supplied dated curation derived from the OWR record, not an independently verified literature assessment. The supplied background explicitly calls for independent verification of time-sensitive status and cited claims.
\end{itemize}
\end{document}
